\section{Related Work}
\label{sec:related_work}

LLM-as-a-judge research studies the reliability, calibration, and agreement of
automatic preference evaluators, including broad judge benchmarks
\citep{zheng2023judging,tan2025judgebench,zhou2025jetts,yang2025reliablejudge}.
Multimodal judge and reward-model work extends that setting to
image-conditioned outputs and preference data
\citep{chen2024mllmjudge,chen2025mjbench,ko2025flexjudge,yasunaga2025multimodal,
hu2026multimodalrewardbench2,zhang2026r1reward,jin2026omnireward}.
These evaluations primarily characterize performance on fixed examples; they
do not ask whether judges with closely matched scores respond similarly to
controlled changes in visual evidence.

Visual spurious-correlation and counterfactual studies test whether visual
models use task-relevant evidence. MM-JudgeBias and Perceptual Judgment Bias
study controlled perturbations in multimodal judges and perceptual bias
\citep{li2025devil,yang2025spuriverse,gao2025scope,lee2026mmjudgebias,
park2026perceptual}. These studies motivate testing behavior under controlled
changes rather than relying on a fixed-example score alone.

VisualFLIP already holds a question fixed, changes task-critical visual
evidence, and evaluates prediction updating \citep{zhu2026visualflip}. Those
ingredients are therefore not themselves contributions of RewardLens.
RewardLens instead asks whether residual intervention-behavior variation
remains among judges matched on an independent static preference score, using
both relevant and irrelevant edit branches with a fixed candidate pair.
Appendix~\ref{app:related-work-scope} makes this boundary explicit.
